% !TeX root = ../../Book3.tex
% 纲要标题：### 16.4 Serre 条件
% 纲要位置：工作记录/计划纲要.md，第 2507 行。
% 纲要细目（写作范围）：
% * 模的 S_k 条件
% * 局部化后的深度要求
% * S_0、S_1、S_2 的层级与嵌入分量的联系
% * 一般 S_k 正则截面的下降；极大点数值不足与非约化反例
% ---

\section{Serre 条件}
\label{b3:sec:1604}

深度在一个点的值不能代替所有局部化处的检查。Serre 条件按各素点的局部维数提出统一下界，使闭点附近的退化与低余维处的性质能够分别比较。


% 纲要内容（仅作写作计划，尚非教材正文）：
%
% * 模的 S_k 条件
% * 局部化后的深度要求
% * S_0、S_1、S_2 的层级与嵌入分量的联系
%
% ---
%

\subsection{模的 \texorpdfstring{\(S_k\)}{Sk} 条件}
\label{b3:subsec:1604:01}

\begin{definition}\label{b3:def:serre-sk}
设 \(R\) 为交换 Noether 环，\(M\) 为有限生成的 \(R\) 模，\(k\ge0\) 为整数。若对每个 \(\mathfrak p\in\operatorname{Supp}_RM\) 都有
\[
\operatorname{depth}_{R_{\mathfrak p}}M_{\mathfrak p}
\ge\min\{k,\dim_{R_{\mathfrak p}}M_{\mathfrak p}\},
\]
则称 \(M\) 满足\term[serre-tiaojian]{Serre 条件}[Serre's condition] \((S_k)\)。对 \(M=R\) 采用同一定义，称环 \(R\) 满足 \((S_k)\)。零模的条件为空。
\end{definition}
\symidx{\((S_k)\)：Serre 深度条件}

深度不超过局部模的维数，所以在 \(\dim M_{\mathfrak p}\le k\) 的点，\((S_k)\) 要求深度恰等于该维数；在局部维数至少为 \(k\) 的点，则只要求深度至少为 \(k\)。维数为零的点没有额外要求。

这里采用局部模的维数，不能把它统一换成 \(\dim R_{\mathfrak p}\)：模的支集可能只有环境的一部分。例如，设 \(k\) 为域，\(R=k[X,Y]_{(X,Y)}\) 上的 \(M=R/(X)\) 只支撑在直线 \(X=0\) 上。在闭点，\(M\cong k[Y]_{(Y)}\) 的深度与模维数都为一；在这条线的一般点，局部模为 \(k(Y)\)，深度与维数都为零。因此它满足全部 Serre 条件，而不必具有深度二。

\subsection{局部化后的深度要求}
\label{b3:subsec:1604:02}

\begin{proposition}\label{b3:prop:serre-localization}
设 \(R\) 为交换 Noether 环，\(M\) 为有限生成的 \(R\) 模，\(k\ge0\) 为整数，\(S\subseteq R\) 为乘法集。
\begin{enumerate}
\item 若 \(M\) 满足 \((S_k)\)，则 \(S^{-1}M\) 在 \(S^{-1}R\) 上满足 \((S_k)\)。
\item \(M\) 满足 \((S_k)\) 当且仅当每个极大理想处的 \(M_{\mathfrak m}\) 满足 \((S_k)\)。
\item \((S_{k+1})\) 蕴含 \((S_k)\)，而每个有限生成模都满足 \((S_0)\)。
\end{enumerate}
\end{proposition}
\begin{proof}
局部化环的素理想对应于原环中避开乘法集的素理想，且在这些素理想处再局部化所得的环和模与直接局部化相同。因此深度与模维数都相同，第一条成立。任意素理想包含于某个极大理想，在该极大理想处先局部化再取原素理想，对象仍相同，得到第二条的反向；正向来自第一条。第三条直接比较最小值，局部环上非零有限生成模的深度非负。
\end{proof}

注意第二条要求极大理想处的局部模本身满足 \((S_k)\)，因而已经检查其全部素理想处的条件。只检查原环各极大理想处的一条深度数值，并不能自动完成定义中的所有检查。

\subsection{\texorpdfstring{\(S_1\)}{S1}、\texorpdfstring{\(S_2\)}{S2} 与嵌入分量}
\label{b3:subsec:1604:03}

\begin{theorem}\label{b3:thm:serre-s1-associated}
设 \(R\) 为交换 Noether 环，\(M\ne0\) 为有限生成的 \(R\) 模。则 \(M\) 满足 \((S_1)\) 当且仅当
\[
\operatorname{Ass}_RM=\operatorname{Min}(\operatorname{Supp}_RM),
\]
即其关联素理想全是支集的极小素理想。
\end{theorem}
\begin{proof}
对每个 \(\mathfrak p\in\operatorname{Supp}M\)，\(M_{\mathfrak p}\) 是非零有限生成的 \(R_{\mathfrak p}\) 模。由\cref{b3:cor:depth-zero-socle}，其深度为零当且仅当局部极大理想 \(\mathfrak pR_{\mathfrak p}\) 是关联素理想；再用\cref{b3:prop:associated-localization}及素理想的收缩对应，这又等价于 \(\mathfrak p\in\operatorname{Ass}_RM\)。另一方面，\(\dim M_{\mathfrak p}=0\) 当且仅当支集中没有严格小于 \(\mathfrak p\) 的素理想，即 \(\mathfrak p\) 是支集的极小元。

\((S_1)\) 只要求局部维数为正的各点深度不为零。因此它恰好排除非极小的关联素理想。由\cref{b3:prop:minimal-primes-associated}，支集的极小素理想总是关联素理想，所以排除这些嵌入关联素理想就得到所列集合等号。
\end{proof}

对域 \(k\)，第四章的两轴环 \(k[x,y]/(xy)\) 只有由 \((x)\)、\((y)\) 给出的两个极小关联素理想，所以满足 \((S_1)\)。而 \(k[x,y]/(x^2,xy)\) 中非零的 \(\bar x\) 被 \((x,y)\) 零化，其零化子恰为这个极大理想；它严格包含极小素理想 \((x)\)，所以该环不满足 \((S_1)\)。在这个闭点，局部维数仍为一，深度却已经降为零。

因此，对交换 Noether 环 \(R\) 上的有限生成模 \(M\)，这三层条件可以写成
\[
\begin{aligned}
(S_0)&:\quad\text{自动成立},\\
(S_1)&:\quad\text{没有嵌入关联素理想},\\
(S_2)&:\quad(S_1)\text{，并且在}\;\dim M_{\mathfrak p}\ge2
\;\text{的各点有}\;\operatorname{depth}M_{\mathfrak p}\ge2.
\end{aligned}
\]
它们满足 \((S_2)\Rightarrow(S_1)\Rightarrow(S_0)\)。这里 \((S_2)\) 仍包含一维局部模的深度至少为一的要求，不能只检查二维及更高维的点。

\ProseParagraphBreak

\((S_1)\) 与 \((S_2)\) 的差别可以在一个简单的理想上看见。设 \(k\) 为域，取 \(R=k[X,Y]_{(X,Y)}\)、\(M=(X,Y)\)。\(M\) 为整环中的非零理想，所有非零元的零化子都是零，所以它无挠，\(\operatorname{Ass}M=\{0\}\)，且支集为整个 \(\operatorname{Spec}R\)。它满足 \((S_1)\)；但前节算出其深度为一，模维数为二，故不满足 \((S_2)\)。排除嵌入关联素理想只保证正维局部点有第一层非零因子，并不保证二维及更高维处还能继续正则取商。

\begin{proposition}\label{b3:prop:s2-regular-section-s1}
设 \(R\) 为交换 Noether 环，\(s\ge1\) 为整数，\(M\) 为满足 \((S_s)\) 的有限生成的 \(R\) 模，\(x\in R\) 在 \(M\) 上为非零因子。则 \(M/xM\) 满足 \((S_{s-1})\)。特别地，满足 \((S_2)\) 的模沿非零因子取商后满足 \((S_1)\)。
\end{proposition}
\begin{proof}
局部化保持 \(x\) 的乘法单射。若 \(M_{\mathfrak p}\ne0\) 且 \(x\in\mathfrak p\)，Nakayama 引理保证 \(M_{\mathfrak p}/xM_{\mathfrak p}\ne0\)；若 \(x\notin\mathfrak p\)，则 \(x\) 成为单位，商模的局部化为零。因此
\[
\operatorname{Supp}(M/xM)=\operatorname{Supp}M\cap V(x).
\]
取这个支集中的 \(\mathfrak p\)，\(x/1\) 就是非零有限生成的 \(R_{\mathfrak p}\) 模 \(M_{\mathfrak p}\) 上的正则元素。记 \(d=\dim M_{\mathfrak p}\)、\(t=\operatorname{depth}M_{\mathfrak p}\)，由\cref{b3:prop:regular-quotient-dimension,b3:cor:regular-quotient-depth}，商的局部维数为 \(d-1\)，深度为 \(t-1\)。其中 \(d\ge1\)，而 \((S_s)\) 给出 \(t\ge\min\{s,d\}\)，故
\[
t-1\ge\min\{s,d\}-1=\min\{s-1,d-1\}.
\]
这正是商模在 \(\mathfrak p\) 处所需的 \((S_{s-1})\) 不等式。若商模为零，条件为空，结论也成立。
\end{proof}

因此 \((S_2)\) 不只排除当前模的嵌入部分，也控制正则截面不会新生这种部分。它仍不是正规性的完整条件；\chapref{b3:ch:20}将另外加入高度不超过一的各素点处的正则性，并证明两种要求如何共同刻画正规环。

\begin{proposition}\label{b3:prop:serre-condition-counterexamples}
非零交换 Artin 环满足所有 Serre 条件 \((S_s)\)，其中 \(s\ge0\) 为整数，但不必约化。例如，对任意域 \(k\)，环 \(k[\varepsilon]/(\varepsilon^2)\) 满足全部这些条件而不约化。

另设 \(k\) 为域，\(R=k[X,Y,Z]_{(X,Y,Z)}\)、\(\mathfrak m=(X,Y,Z)R\)、\(M=R\oplus R/(X,Y)\)。则 \(\operatorname{depth}_RM=1\)，但在素理想 \(\mathfrak p=(X,Y)R\) 处，有
\[
\operatorname{depth}_{R_{\mathfrak p}}M_{\mathfrak p}=0,
\qquad \dim_{R_{\mathfrak p}}M_{\mathfrak p}=2.
\]
因此 \(M\) 满足闭点自身的 \((S_1)\) 深度下界，却不满足 \((S_1)\)。
\end{proposition}
\begin{proof}
交换 Artin 环由\cref{b3:thm:artin-finite-length}是 Noether 环，且由\cref{b3:lem:artin-primes-radical}，每个素理想都极大。因此在每个素点，局部环的维数为零，Serre 条件只要求深度非负，自动成立。在 \(k[\varepsilon]/(\varepsilon^2)\) 中，\(\varepsilon\ne0\) 而 \(\varepsilon^2=0\)，所以约化性不能由这些深度条件推出。

在所列 \(R\) 与 \(M\) 中，\(R/(X,Y)\cong k[Z]_{(Z)}\)，所以 \(Z\) 在两个直和分量上都为非零因子，也就在 \(M\) 上为非零因子。商模
\[
M/ZM\cong R/(Z)\oplus k
\]
含有非零元素 \((0,1)\)，其零化子为 \(\mathfrak m\)。由\cref{b3:cor:depth-zero-socle}，该商模深度为零；再由\cref{b3:cor:regular-quotient-depth}，\(\operatorname{depth}_RM=1\)。

在 \(\mathfrak p=(X,Y)R\) 处，局部化得到
\[
R_{\mathfrak p}\cong k(Z)[X,Y]_{(X,Y)},\qquad
M_{\mathfrak p}\cong R_{\mathfrak p}\oplus k(Z).
\]
第二个分量成为局部环的剩余域，故同样有非零元素被整个局部极大理想零化，给出深度零。第一个分量是环本身，所以局部模的支集为整个局部谱。这个局部环有素链 \((0)\subsetneq(X)\subsetneq(X,Y)\)，而极大理想由两个元素生成，高度定理给出其维数恰为二。因此在这个支集点，\((S_1)\) 要求深度至少为一，实际却为零。
\end{proof}
